CUDA FortranでもOpen ACCのようなdirectiveベースの並列化を実装できる．コンパイラの判断でimplicit private等の属性を付与してくれる．Reduction処理を扱うこともできる．コンパイラオプションに-Minfo=accelを付与することでコンパイラがどのように並列化したかを確認することができる．

OUxSBLIでは保存量から基本量を計算する処理にdirectiveベースの並列化を用いている．以下にコードを示す．

\begin{lstlisting}[language=Fortran]
!$cuf kernel do(3)<<<*,(32,4,2)>>>
do k = 1, nz
  do j = 1, ny
    do i = 1, nx
      over_Q1 = 1.d0 / QJ(1,i,j,k)
      rho     = Jacobian(i,j) * QJ(1,i,j,k)
      u       = QJ(2,i,j,k) * over_Q1
      v       = QJ(2,i,j,k) * over_Q1
      w       = QJ(2,i,j,k) * over_Q1
      p       = gamma_1 * (Jacobian(i,j) * QJ(5,i,j,k) - 0.5d0 * rho * (u*u + v*v + w*w))
      Q(1,i,j,k) = rho
      Q(2,i,j,k) = u
      Q(3,i,j,k) = v
      Q(4,i,j,k) = w
      Q(5,i,j,k) = p
      T(i,j,k)   = p / (R * rho)
enddo;enddo;enddo
\end{lstlisting}

\noindent
Directiveベースの並列化ではコンパイラは”安全”な並列化を行うため，本質的にカーネルを起動する並列化と同じコードを書いたとしてもレジスタの使用量等に違いがでる場合がある．また，GPUは割り算がとにかく苦手なので，事前に逆数を計算するのがよい．
